\documentclass[11pt,a4paper]{article}
\usepackage[T2A]{fontenc}
\usepackage[utf8]{inputenc}
\usepackage[russian]{babel}
\usepackage{amsmath,amssymb,amsfonts}
\usepackage{geometry}
\usepackage{array}
\usepackage{booktabs}
\usepackage{longtable}
\usepackage{pdflscape}
\usepackage{enumitem}
\usepackage{hyperref}
\usepackage[table]{xcolor}
\geometry{margin=1.8cm}
\setlength{\parindent}{0pt}
\setlength{\parskip}{0.45em}
\newcommand{\sgop}{\mathop{\sum\!\!\!\int_{\bowtie}}}

\title{IIS Full Test + Documentation Pack}
\author{Проект анализа интегрального индекса состояния}
\date{\today}

\begin{document}
\maketitle

\section{Введение}
Этот отчёт объединяет в одном документе версии \texttt{IISVersion1}--\texttt{IISVersion7}, полную матрицу доступных тестовых сценариев, state capacity, Dreamer-калибровки и synthetic-бенч оператора Sumigron.

Главный принцип документа: слабые результаты не скрываются. Если версия не универсальна, если режим не адаптирован к динамике или если на конкретном датасете видны только tentative-сигналы, это фиксируется явно.

\textbf{Лучшая практическая линия сейчас:} IISVersion6.\\
\textbf{Лучшая теоретическая линия:} IISVersion7 + Sumigron.\\
\textbf{Самый сильный реальный сценарий:} ds002724 / strict (V6).

\section{Что такое strict / proxy / hybrid}
\begin{center}
\begin{tabular}{p{2.2cm}p{11.8cm}}
\toprule
Режим & Смысл \\
\midrule
Strict & Используются только реально измеренные прямые признаки. Ничего не додумывается и не якорится proxy-метками. \\
Proxy & Разрешены proxy-замены и мягкие якоря по доступным косвенным признакам и self-report. Это полезно для исторических V1–V3 и части ранних экспериментов. \\
Hybrid & Модель использует всё, что реально доступно: прямые признаки остаются прямыми, недостающие части могут компенсироваться proxy-слоем. \\
\bottomrule
\end{tabular}
\end{center}

\section{Карта версий и схемы}
\subsubsection*{Эволюция версий V1\ensuremath{\rightarrow}V7} Переход от исторических proxy-ориентированных формул к негормональным и затем к gated/Sumigron-архитектурам.
\begin{center}\fbox{\parbox{2.7cm}{\centering V1: Тестовое раннее ядро}} $\rightarrow$ \fbox{\parbox{2.7cm}{\centering V2: Версия 1.0 с общей лево-правой мощностью}} $\rightarrow$ \fbox{\parbox{2.7cm}{\centering V3: Версия 2.0 с K и отдельной валентностью Q}} $\rightarrow$ \fbox{\parbox{2.7cm}{\centering V4: Первое негормональное практическое ядро}}\\[0.6em]
$\downarrow$\\[0.6em]
\fbox{\parbox{2.7cm}{\centering V5: Контрастная нелинейная версия}} $\rightarrow$ \fbox{\parbox{2.7cm}{\centering V6: Gated-модель режимов}} $\rightarrow$ \fbox{\parbox{2.7cm}{\centering V7: Экспериментальная Sumigron-бета}}\end{center}
\subsubsection*{Общий pipeline} Общая цепочка обработки для всех современных версий.
\begin{center}\fbox{\parbox{2.7cm}{\centering dataset}} $\rightarrow$ \fbox{\parbox{2.7cm}{\centering segments / windows}} $\rightarrow$ \fbox{\parbox{2.7cm}{\centering features}} $\rightarrow$ \fbox{\parbox{2.7cm}{\centering components}} $\rightarrow$ \fbox{\parbox{2.7cm}{\centering IIS}}\end{center}
\subsubsection*{Режимы данных} Три режима доступа к данным и компенсации пропусков.
\begin{center}\fbox{\parbox{2.7cm}{\centering strict}} $\rightarrow$ \fbox{\parbox{2.7cm}{\centering proxy}} $\rightarrow$ \fbox{\parbox{2.7cm}{\centering hybrid}}\end{center}
\subsubsection*{Сравнение архитектур V4/V5/V6/V7} Упрощённая карта того, какие архитектурные идеи добавлялись поверх практического ядра.
\begin{center}\begin{tabular}{lccccc}\toprule Версия & Гормоны & Нелинейная геометрия & Gated-режимы & 3D слой & Sumigron \\\midrule V4 & нет & частично & нет & частично & нет \\ V5 & нет & да & нет & частично & нет \\ V6 & нет & да & да & частично & нет \\ V7 beta & нет & да & да & да & да \\\bottomrule\end{tabular}\end{center}
\subsubsection*{Архитектура V7 beta} Новая ветка: структурный оператор до gated-слоя и 3D-интерпретации.
\begin{center}\fbox{\parbox{2.7cm}{\centering window / subterms}} $\rightarrow$ \fbox{\parbox{2.7cm}{\centering Sumigron}} $\rightarrow$ \fbox{\parbox{2.7cm}{\centering A / Gamma / V / Q}} $\rightarrow$ \fbox{\parbox{2.7cm}{\centering gated IIS}} $\rightarrow$ \fbox{\parbox{2.7cm}{\centering RES / DYN}}\end{center}
\subsubsection*{Внутренняя логика Sumigron} Разделение окна на уровень, структуру и энергию с attentive-взвешиванием.
\begin{center}\fbox{\parbox{2.7cm}{\centering уровень}} $\rightarrow$ \fbox{\parbox{2.7cm}{\centering структура}} $\rightarrow$ \fbox{\parbox{2.7cm}{\centering энергия}} $\rightarrow$ \fbox{\parbox{2.7cm}{\centering attentive weighting}} $\rightarrow$ \fbox{\parbox{2.7cm}{\centering структурный скаляр}}\end{center}
\subsubsection*{Матрица применимости версий} Явное разделение статической, динамической и proxy/hybrid-применимости.
\begin{center}\begin{tabular}{lccccc}\toprule Версия & static & dynamic & proxy & hybrid & calibrated \\\midrule V1 & да & нет & да & да & нет \\ V2 & да & нет & да & да & нет \\ V3 & да & нет & да & да & нет \\ V4 & да & да & да & да & нет \\ V5 & да & да & да & да & нет \\ V6 & да & да & да & да & нет \\ V7 & да & да & да & да & да \\\bottomrule\end{tabular}\end{center}

\section{Полная документация по каждой версии}

\subsection*{V1. Тестовое раннее ядро \textcolor{green!20!black}{\textbf{рабочее}}}
\[
IIS_1=\sigma(0.35A+0.25\Gamma+0.30H+0.10V)
\]
\textbf{Назначение.} Проверить саму идею интегрального индекса на четырёх базовых блоках.

\textbf{Сильная сторона.} Максимально простая и прозрачная структура.

\textbf{Слабое место.} Сильно зависит от гормонального proxy-блока и плохо переносится в strict-режим.

\textbf{Уровень универсальности.} Низкая

\textbf{Статус.} Историческая proxy-ориентированная версия. Полезна как точка отсчёта, но не как современное универсальное ядро.

\textbf{Входные компоненты:}
\begin{itemize}[leftmargin=1.2cm]
\item A по alpha-асимметрии
\item Gamma
\item H
\item V
\end{itemize}

\textbf{Ограничения:}
\begin{itemize}[leftmargin=1.2cm]
\item историческая версия
\item не универсальна
\item proxy-ориентирована
\item не адаптирована к динамике
\end{itemize}


\subsection*{V2. Версия 1.0 с общей лево-правой мощностью \textcolor{green!20!black}{\textbf{рабочее}}}
\[
IIS_2=\sigma(0.35A+0.25\Gamma+0.30H+0.10V)
\]
\textbf{Назначение.} Сохранить простую интеграцию V1, но заменить A на более грубую межполушарную мощность.

\textbf{Сильная сторона.} Чуть ближе к общим EEG-признакам, чем V1.

\textbf{Слабое место.} Остаётся завязанной на гормональный блок и в реальных strict-сценариях почти неработоспособна.

\textbf{Уровень универсальности.} Низкая

\textbf{Статус.} Историческая промежуточная версия. Применимость ограничена proxy/hybrid-сценариями.

\textbf{Входные компоненты:}
\begin{itemize}[leftmargin=1.2cm]
\item A по общей мощности полушарий
\item Gamma
\item H
\item V
\end{itemize}

\textbf{Ограничения:}
\begin{itemize}[leftmargin=1.2cm]
\item историческая версия
\item не универсальна
\item proxy-ориентирована
\item не адаптирована к динамике
\end{itemize}


\subsection*{V3. Версия 2.0 с K и отдельной валентностью Q \textcolor{green!20!black}{\textbf{рабочее}}}
\[
IIS_3=\sigma(0.35A+0.25\Gamma+0.30H+0.10V),\quad H=\frac{K\cdot DA}{50}-\frac{C}{15}
\]
\textbf{Назначение.} Добавить коэффициент проницаемости K и отделить диагностическую валентность Q от основного IIS.

\textbf{Сильная сторона.} Содержательно богаче V1/V2 за счёт K и Q.

\textbf{Слабое место.} Слишком тяжело опирается на proxy-нормировки и не подтверждается как универсальная практическая модель.

\textbf{Уровень универсальности.} Низкая

\textbf{Статус.} Историческая расширенная версия. Полезна для истории идеи и сопоставления, но не для основной практической линии.

\textbf{Входные компоненты:}
\begin{itemize}[leftmargin=1.2cm]
\item A
\item Gamma
\item H через K
\item V
\item Q как отдельный диагностический блок
\end{itemize}

\textbf{Ограничения:}
\begin{itemize}[leftmargin=1.2cm]
\item историческая версия
\item не универсальна
\item proxy-ориентирована
\item не адаптирована к динамике
\end{itemize}


\subsection*{V4. Первое негормональное практическое ядро \textcolor{green!20!black}{\textbf{рабочее}}}
\[
IIS_4=\sigma(0.10A-0.05\Gamma+0.25V+0.60Q_4)
\]
\textbf{Назначение.} Убрать гормоны и собрать рабочую EEG/ECG-модель для открытых данных.

\textbf{Сильная сторона.} Первая реально переносимая строгая версия без выдуманных биохимических блоков.

\textbf{Слабое место.} Шкала ещё довольно плоская; промежуточные режимы часто сглаживаются.

\textbf{Уровень универсальности.} Средняя

\textbf{Статус.} Рабочая практическая база. Хороша как baseline, но по внутренней геометрии уступает V5/V6.

\textbf{Входные компоненты:}
\begin{itemize}[leftmargin=1.2cm]
\item A по alpha/total asymmetry
\item Gamma по gamma и gamma/alpha
\item V по HR и HF/LF
\item Q как физиологическая интеграция
\end{itemize}

\textbf{Ограничения:}
\begin{itemize}[leftmargin=1.2cm]
\item явных специальных флагов нет
\end{itemize}


\subsection*{V5. Контрастная нелинейная версия \textcolor{green!20!black}{\textbf{рабочее}}}
\[
IIS_5=\operatorname{clip}\!\left((1-m)IIS_{5,\mathrm{base}}+mIIS_{5,\mathrm{contrast}},0,1\right)
\]
\textbf{Назначение.} Раскрыть контраст режимов и уменьшить плоскость середины шкалы без возврата к гормонам.

\textbf{Сильная сторона.} Часто лучше выделяет скрытую структуру и средние режимы, чем V4.

\textbf{Слабое место.} Не универсальна по датасетам: на части наборов выигрывает, на части нет; чувствительна к калибровке.

\textbf{Уровень универсальности.} Средняя

\textbf{Статус.} Нелинейная практическая версия. Хорошо работает там, где помогает контрастная геометрия, но не является окончательно универсальной.

\textbf{Входные компоненты:}
\begin{itemize}[leftmargin=1.2cm]
\item A
\item Gamma
\item V
\item Q с coherence/synergy/conflict
\end{itemize}

\textbf{Ограничения:}
\begin{itemize}[leftmargin=1.2cm]
\item не универсальна
\end{itemize}


\subsection*{V6. Gated-модель режимов \textcolor{green!20!black}{\textbf{рабочее}}}
\[
IIS_6=g_{\mathrm{reg}}i_{\mathrm{reg}}+g_{\mathrm{mob}}i_{\mathrm{mob}}+g_{\mathrm{dep}}i_{\mathrm{dep}}-\lambda_H H(g)-\lambda_C C
\]
\textbf{Назначение.} Сделать модель смесью режимов regulation / mobilization / depletion вместо одной глобальной почти линейной формулы.

\textbf{Сильная сторона.} Лучшее текущее универсальное практическое strict-ядро, особенно на ds002724.

\textbf{Слабое место.} Число базовых режимов зашито в gate-архитектуру; не всегда раскрывает более богатую карту состояний.

\textbf{Уровень универсальности.} Высокая

\textbf{Статус.} Текущая лучшая практическая линия для честных негормональных strict-сценариев.

\textbf{Входные компоненты:}
\begin{itemize}[leftmargin=1.2cm]
\item A
\item Gamma
\item V
\item Q
\item softmax-gates режимов
\end{itemize}

\textbf{Ограничения:}
\begin{itemize}[leftmargin=1.2cm]
\item явных специальных флагов нет
\end{itemize}


\subsection*{V7. Экспериментальная Sumigron-бета \textcolor{green!20!black}{\textbf{рабочее}}}
\[
IIS_7=g^{reg}_7 s^{reg}_7 + g^{mob}_7 s^{mob}_7 + g^{dep}_7 s^{dep}_7,\quad \sgop=\sum+\int+\bowtie
\]
\textbf{Назначение.} Перенести структурную нелинейность раньше по цепочке через оператор Sumigron и сохранить 3D-слой IIS/RES/DYN.

\textbf{Сильная сторона.} Даёт новую архитектурную линию и уже показывает пользу на части наборов и synthetic-тестах.

\textbf{Слабое место.} Пока beta-версия; есть dataset-specific override для DREAMER, поэтому она ещё не универсальна.

\textbf{Уровень универсальности.} Экспериментальная

\textbf{Статус.} Экспериментальная V7-beta. Есть Dreamer-specific override, поэтому её нельзя пока подавать как окончательно универсальную.

\textbf{Входные компоненты:}
\begin{itemize}[leftmargin=1.2cm]
\item A, Gamma, V через Sumigron-подтермы
\item Q как структурная интеграция
\item gated IIS
\item RES
\item DYN
\end{itemize}

\textbf{Ограничения:}
\begin{itemize}[leftmargin=1.2cm]
\item не универсальна
\item dataset-specific calibration
\end{itemize}


\section{Матрица совместимости}
\begin{center}
\begin{tabular}{lccccc}
\toprule
Версия & static & dynamic & proxy & hybrid & Флаги \\
\midrule
V1 & да & нет & да & да & историческая версия, не универсальна, proxy-ориентирована, не адаптирована к динамике \\
V2 & да & нет & да & да & историческая версия, не универсальна, proxy-ориентирована, не адаптирована к динамике \\
V3 & да & нет & да & да & историческая версия, не универсальна, proxy-ориентирована, не адаптирована к динамике \\
V4 & да & да & да & да & --- \\
V5 & да & да & да & да & не универсальна \\
V6 & да & да & да & да & --- \\
V7 & да & да & да & да & не универсальна, dataset-specific calibration \\
\bottomrule
\end{tabular}
\end{center}

\section{Сравнение версий на тестах}
\subsection*{Полная матрица по всем dataset/mode}
\begin{landscape}
\scriptsize
\begin{longtable}{p{2.8cm}p{0.8cm}rrrrrp{1.2cm}rp{4.4cm}p{1.6cm}}
\toprule
Scenario & V & Cov & Valid & Eff & Ovl & Sens & Rel & k & Verdict & Dyn \\
\midrule
\endfirsthead
\toprule
Scenario & V & Cov & Valid & Eff & Ovl & Sens & Rel & k & Verdict & Dyn \\
\midrule
\endhead
CASE / hybrid & V1 & 1.000 & 7230 & 0.031 & 0.999 & 0.000 & low & 2 & tentative evidence for 2 states & not\_adapted \\
CASE / hybrid & V2 & 1.000 & 7230 & 0.031 & 0.999 & 0.000 & low & 2 & tentative evidence for 2 states & not\_adapted \\
CASE / hybrid & V3 & 1.000 & 7230 & 0.031 & 0.999 & 0.000 & low & 2 & tentative evidence for 2 states & not\_adapted \\
CASE / hybrid & V4 & 1.000 & 7230 & -0.087 & 0.918 & 0.073 & medium & 2 & tentative evidence for 2 states & available \\
CASE / hybrid & V5 & 1.000 & 7230 & -0.097 & 0.916 & 0.152 & medium & 2 & tentative evidence for 2 states & partial\_optional \\
CASE / hybrid & V6 & 1.000 & 7230 & -0.117 & 0.917 & 0.332 & medium & 3 & tentative evidence for 2 states & available \\
CASE / hybrid & V7 & 1.000 & 7230 & -0.040 & 0.905 & 0.022 & medium & 2 & tentative evidence for 2 states & partial\_optional \\
WESAD / hybrid & V1 & 1.000 & 1087 & -0.406 & 0.807 & 0.202 & low & 2 & confidently 2 states & not\_adapted \\
WESAD / hybrid & V2 & 1.000 & 1087 & -0.406 & 0.807 & 0.202 & low & 2 & confidently 2 states & not\_adapted \\
WESAD / hybrid & V3 & 1.000 & 1087 & -0.406 & 0.807 & 0.202 & low & 2 & confidently 2 states & not\_adapted \\
WESAD / hybrid & V4 & 1.000 & 1087 & -0.701 & 0.674 & 0.073 & medium & 2 & confidently 2 states & available \\
WESAD / hybrid & V5 & 1.000 & 1087 & -0.699 & 0.666 & 0.105 & medium & 3 & confidently 2 states, limited evidence for 3rd mode & available \\
WESAD / hybrid & V6 & 1.000 & 1087 & -0.676 & 0.676 & 0.215 & medium & 2 & confidently 2 states & available \\
WESAD / hybrid & V7 & 1.000 & 1087 & -0.694 & 0.648 & 0.061 & medium & 2 & confidently 2 states & partial\_optional \\
ds002722 / strict & V1 & 0.000 & 0 & --- & --- & --- & low & --- & nan & not\_adapted \\
ds002722 / strict & V2 & 0.000 & 0 & --- & --- & --- & low & --- & nan & not\_adapted \\
ds002722 / strict & V3 & 0.000 & 0 & --- & --- & --- & low & --- & nan & not\_adapted \\
ds002722 / strict & V4 & 0.998 & 6756 & 0.048 & 0.867 & 0.033 & high & 3 & tentative evidence for 2 states & available \\
ds002722 / strict & V5 & 0.998 & 6756 & 0.001 & 0.828 & 0.080 & high & 2 & tentative evidence for 2 states & available \\
ds002722 / strict & V6 & 0.998 & 6756 & 0.030 & 0.789 & 0.157 & high & 3 & tentative evidence for 2 states & available \\
ds002722 / strict & V7 & 0.998 & 6756 & 0.000 & 0.797 & 0.084 & high & 2 & tentative evidence for 2 states & available \\
ds002722 / proxy & V1 & 0.999 & 1691 & -2.410 & 0.212 & 0.057 & high & 2 & confidently 2 states, limited evidence for 3rd mode & not\_adapted \\
ds002722 / proxy & V2 & 0.999 & 1691 & -2.412 & 0.224 & 0.057 & high & 2 & confidently 2 states, limited evidence for 3rd mode & not\_adapted \\
ds002722 / proxy & V3 & 0.999 & 1691 & -1.853 & 0.329 & 0.060 & high & 2 & confidently 2 states, limited evidence for 3rd mode & not\_adapted \\
ds002722 / proxy & V4 & 1.000 & 1692 & -0.959 & 0.522 & 0.033 & medium & 3 & likely 3 states & available \\
ds002722 / proxy & V5 & 1.000 & 1692 & -0.290 & 0.675 & 0.094 & medium & 2 & tentative evidence for 2 states & partial\_optional \\
ds002722 / proxy & V6 & 1.000 & 1692 & 0.034 & 0.668 & 0.162 & medium & 3 & tentative evidence for 2 states & available \\
ds002722 / proxy & V7 & 1.000 & 1692 & 0.053 & 0.627 & 0.086 & medium & 2 & tentative evidence for 2 states & partial\_optional \\
ds002724 / strict & V1 & 0.000 & 0 & --- & --- & --- & low & --- & nan & not\_adapted \\
ds002724 / strict & V2 & 0.000 & 0 & --- & --- & --- & low & --- & nan & not\_adapted \\
ds002724 / strict & V3 & 0.000 & 0 & --- & --- & --- & low & --- & nan & not\_adapted \\
ds002724 / strict & V4 & 0.998 & 13801 & -0.371 & 0.825 & 0.033 & high & 2 & confidently 2 states & available \\
ds002724 / strict & V5 & 0.998 & 13801 & -0.413 & 0.802 & 0.072 & high & 2 & confidently 2 states & available \\
ds002724 / strict & V6 & 0.998 & 13801 & -0.433 & 0.799 & 0.144 & high & 2 & confidently 2 states & available \\
ds002724 / strict & V7 & 0.998 & 13801 & -0.254 & 0.845 & 0.067 & high & 2 & tentative evidence for 2 states & available \\
ds002724 / proxy & V1 & 0.999 & 3454 & -2.687 & 0.174 & 0.058 & high & 2 & confidently 2 states, limited evidence for 3rd mode & not\_adapted \\
ds002724 / proxy & V2 & 0.999 & 3454 & -2.693 & 0.174 & 0.058 & high & 3 & likely 3 states & not\_adapted \\
ds002724 / proxy & V3 & 0.999 & 3454 & -2.178 & 0.273 & 0.061 & high & 4 & likely 3 states & not\_adapted \\
ds002724 / proxy & V4 & 1.000 & 3456 & -1.446 & 0.455 & 0.033 & medium & 2 & confidently 2 states, limited evidence for 3rd mode & available \\
ds002724 / proxy & V5 & 1.000 & 3456 & -0.779 & 0.602 & 0.079 & medium & 3 & confidently 2 states, limited evidence for 3rd mode & partial\_optional \\
ds002724 / proxy & V6 & 1.000 & 3456 & -0.564 & 0.689 & 0.145 & medium & 3 & confidently 2 states, limited evidence for 3rd mode & available \\
ds002724 / proxy & V7 & 1.000 & 3456 & -0.536 & 0.639 & 0.067 & medium & 3 & confidently 2 states, limited evidence for 3rd mode & partial\_optional \\
DREAMER / strict & V1 & 0.000 & 0 & --- & --- & --- & low & --- & nan & not\_adapted \\
DREAMER / strict & V2 & 0.000 & 0 & --- & --- & --- & low & --- & nan & not\_adapted \\
DREAMER / strict & V3 & 0.000 & 0 & --- & --- & --- & low & --- & nan & not\_adapted \\
DREAMER / strict & V4 & 0.711 & 23169 & 0.025 & 0.886 & 0.033 & high & 2 & tentative evidence for 2 states & available \\
DREAMER / strict & V5 & 0.711 & 23169 & -0.006 & 0.876 & 0.063 & high & 2 & tentative evidence for 2 states & available \\
DREAMER / strict & V6 & 0.711 & 23169 & 0.004 & 0.893 & 0.157 & high & 2 & tentative evidence for 2 states & available \\
DREAMER / strict & V7 & 0.711 & 23169 & -0.016 & 0.897 & 0.061 & high & 3 & tentative evidence for 2 states & available \\
\bottomrule
\end{longtable}
\normalsize
\end{landscape}

\subsection*{Лучшие версии по сценариям}
\begin{center}
\begin{tabular}{p{2.8cm}p{0.8cm}p{1.3cm}rrrrp{4.2cm}}
\toprule
Scenario & V & Rel & Eff & Ovl & Sens & Verdict \\
\midrule
CASE / hybrid & V6 & medium & -0.117 & 0.917 & 0.332 & tentative evidence for 2 states \\
DREAMER / strict & V4 & high & 0.025 & 0.886 & 0.033 & tentative evidence for 2 states \\
ds002722 / proxy & V1 & high & -2.410 & 0.212 & 0.057 & confidently 2 states, limited evidence for 3rd mode \\
ds002722 / strict & V6 & high & 0.030 & 0.789 & 0.157 & tentative evidence for 2 states \\
ds002724 / proxy & V2 & high & -2.693 & 0.174 & 0.058 & likely 3 states \\
ds002724 / strict & V6 & high & -0.433 & 0.799 & 0.144 & confidently 2 states \\
WESAD / hybrid & V6 & medium & -0.676 & 0.676 & 0.215 & confidently 2 states \\
\bottomrule
\end{tabular}
\end{center}

\subsection*{Что реально подтверждается по числу состояний}
\begin{center}
\begin{tabular}{p{2.8cm}p{0.8cm}p{4.2cm}p{1.3cm}p{6.0cm}}
\toprule
Scenario & V & Подтверждённое & Потолок & Примечание \\
\midrule
CASE / hybrid & V6 & tentative evidence for 2 states & k\ensuremath{\approx}3 & Внутренняя геометрия модели дробится минимум на два режима, но размеченные состояния пока не разделяются достаточно уверенно. \\
WESAD / hybrid & V6 & confidently 2 states & k\ensuremath{\approx}2 & Бинарное разделение устойчиво подтверждается, а дополнительные режимы не получают достаточно сильной кластерной и попарной поддержки. \\
ds002722 / strict & V6 & tentative evidence for 2 states & k\ensuremath{\approx}3 & Внутренняя геометрия модели дробится минимум на два режима, но размеченные состояния пока не разделяются достаточно уверенно. \\
ds002722 / proxy & V1 & confidently 2 states, limited evidence for 3rd mode & k\ensuremath{\approx}2 & Есть намёки на переходный или промежуточный режим, но устойчивости недостаточно для честного подтверждения трёх состояний. \\
ds002724 / strict & V6 & confidently 2 states & k\ensuremath{\approx}2 & Бинарное разделение устойчиво подтверждается, а дополнительные режимы не получают достаточно сильной кластерной и попарной поддержки. \\
ds002724 / proxy & V2 & likely 3 states & k\ensuremath{\approx}3 & Третий режим поддерживается и попарной статистикой по меткам, и внутренней геометрией модели, но слабее первых двух. \\
DREAMER / strict & V4 & tentative evidence for 2 states & k\ensuremath{\approx}2 & Внутренняя геометрия модели дробится минимум на два режима, но размеченные состояния пока не разделяются достаточно уверенно. \\
\bottomrule
\end{tabular}
\end{center}

\section{V7 beta и Sumigron}
\textbf{Зачем появился оператор.} Sumigron нужен как попытка переставить нелинейность раньше по цепочке: не сжимать окно в один ранний bounded-ответ, а сначала сохранить уровень, структуру и энергию.

\textbf{Что делает V7 beta.} Архитектура \(\text{window} \to \sgop \to A,\Gamma,V,Q \to IIS,RES,DYN\) сохраняет gated-слой и трёхмерную интерпретацию, но меняет раннюю агрегацию.

\textbf{Где уже дал пользу.} На части строгих сценариев и в synthetic-бенче Sumigron/V7 показывают осмысленный выигрыш.

\textbf{Где пока не дал прорыва.} На DREAMER V7 полезна, но не делает резкого скачка по подтверждённому числу состояний; кроме того, для DREAMER уже используется dataset-specific override, поэтому версия ещё не универсальна.

\section{Экспериментальные проверки}
\subsection*{Dreamer V5 calibration}
Current global objective: 0.342.\\
Best candidate objective: 0.496.\\
Калибровка V5 под DREAMER показывает, что часть выигрыша достигается именно доменной подстройкой нелинейных параметров.

\subsection*{Dreamer V7 beta calibration}
Candidate mode: global.\\
Current global objective: 0.510.\\
Best candidate objective: 0.523.\\
Для V7 beta зафиксирован отдельный Dreamer-specific override; это улучшение локально, но не делает модель глобально универсальной.

\subsection*{Synthetic Sumigron benchmark}
Structure-driven / Sumigron \(\epsilon^2\): 0.635.\\
Structure-driven / sigmoid \(\epsilon^2\): 0.012.\\
Synthetic-бенч показывает, что Sumigron особенно силён там, где состояния различаются не только средним уровнем, но и структурой окна.

\section{Ограничения и честные пометки}
\begin{itemize}[leftmargin=1.2cm]
\item \textbf{V1}: Сильно зависит от гормонального proxy-блока и плохо переносится в strict-режим. Ограничения: историческая версия, не универсальна, proxy-ориентирована, не адаптирована к динамике.
\item \textbf{V2}: Остаётся завязанной на гормональный блок и в реальных strict-сценариях почти неработоспособна. Ограничения: историческая версия, не универсальна, proxy-ориентирована, не адаптирована к динамике.
\item \textbf{V3}: Слишком тяжело опирается на proxy-нормировки и не подтверждается как универсальная практическая модель. Ограничения: историческая версия, не универсальна, proxy-ориентирована, не адаптирована к динамике.
\item \textbf{V4}: Шкала ещё довольно плоская; промежуточные режимы часто сглаживаются. Ограничения: без специальных флагов.
\item \textbf{V5}: Не универсальна по датасетам: на части наборов выигрывает, на части нет; чувствительна к калибровке. Ограничения: не универсальна.
\item \textbf{V6}: Число базовых режимов зашито в gate-архитектуру; не всегда раскрывает более богатую карту состояний. Ограничения: без специальных флагов.
\item \textbf{V7}: Пока beta-версия; есть dataset-specific override для DREAMER, поэтому она ещё не универсальна. Ограничения: не универсальна, dataset-specific calibration.
\end{itemize}

\textbf{DREAMER.} DREAMER полезен для выхода за грубые два режима, но пока остаётся слабым по подтверждённому числу состояний.

\textbf{Sumigron.} Sumigron остаётся экспериментальным оператором: польза уже видна на synthetic и части V7-прогонов, но это ещё не финальная доказанная основа всей системы.

\section{Заключение}
\textbf{Текущая лучшая практическая версия.} V6 --- Gated-модель режимов.

\textbf{Текущая лучшая теоретическая линия.} IISVersion7 + Sumigron.

\textbf{Что остаётся гипотезой.} Подтверждение 3+ устойчивых состояний на разнородных наборах, а также финальная универсальность V7/Sumigron без dataset-specific override.

\end{document}
